\section{Обобщённые координаты}
\label{sec:coords-dh-dt}
\label{ch:native-coords}

\begin{definition}[Внешние координаты \texorpdfstring{$\caCoordsDhDt$}{dh,dt}]
\label{def:dh-dt}
Пусть $\caStepSpace$ --- шаг решётки, $\htick$ --- длительность такта.
\textbf{Внешние координаты} $\caCoordsDhDt$ на~$\layerMicro$ --- это пара
\[
  dh:=\caStepSpace,\qquad dt:=\htick,
\]
при $\caSignalSpeed=\caStepSpace/\htick$.
В натуральных единицах симуляции полагают $dh=dt=1$.
Под \textbf{ступенью действия} $\ladderAction$ на $(\hv,1\,\htick)$ понимают квант, из которого
одной подстановкой следуют $\ladderMomentum=\ladderAction/\caStepSpace$,
$\ladderEnergy=\ladderAction/\htick$ и $\ladderForce=\ladderAction/(\caStepSpace\cdot\htick)$,
без трёх независимых подгонок.
\end{definition}

\begin{definition}[Внутренние координаты]
\label{def:internal-coords}
Пусть регистр элементарной ячейки задаётся $(\symMu,k_{\symPhi},\phaseVarphi_{\mathrm{f}})$
или, эквивалентно, $(\symMu,\phaseVarphi_{\mathrm{disc}})$.
Эту запись и называют \textbf{внутренними координатами}; спинор~$\spaceC^2$ хранит две копии регистра.
\textbf{Вакуум} в этой записи --- $\phaseVarphi_{\mathrm{f}}=0$, $\phaseVarphi_{\mathrm{disc}}=k_{\symPhi}\cdot 41$.
\end{definition}

\begin{table}[ht]
\centering
\caption{Лестница $\ladderAction$ и мост $\layerMicro\to\layerMacro$.}
\label{tab:dh-dt-bridge}
\begin{tabular}{@{}lll@{}}
\toprule
На $\layerMicro$ & $\caCoordsDhDt$ & Мост $\layerMacro$ \\
\midrule
$\ladderAction$ & на $(\hv,1\,\htick)$ & $\planckHbar/2$ \\
$\ladderMomentum$ & $\ladderAction/\caStepSpace$ & $\planckHbar/(2\caStepSpace)$ \\
$\ladderEnergy$ & $\ladderAction/\htick$ & $\planckHbar/(2\htick)$ \\
$\caSignalSpeed$ & $\caStepSpace/\htick$ & $\caStepSpace/\htick$ \\
\bottomrule
\end{tabular}
\end{table}

